\clearpage
\section*{September 22, 2026: Land measurement barely moves the bunching comparison}
\textit{Agent-drafted research entry.}

The parcel-boundary review had flagged about 150 weighting-sample parents for
case-by-case documentary research. Before continuing it, we asked whether land
measurement can change the reweighted comparison at all. The answer is that
it cannot change the bunching findings. The rule used to measure land does
matter for the size of historical projects, and the automatic correction the
audit offers applies unevenly across the two periods.

\paragraph{The test.} The benchmark sample is unchanged: 595
historical and 310 post-period rental parents with at least 50 units.
Each scenario recalibrates the historical weights on the same four moments
(log lot area, residential FAR, built FAR, borough). Units, filings and
membership never change. The scenarios are: no weighting; production land;
dropping every flagged parent from both periods; replacing lot area with the
audit's earlier-parcel area for parents that pass the boundary screen; and
halving or doubling every flagged parent's land while holding earlier
building floor fixed.

\begin{center}
\small
\begin{tabular}{lrrrrr}
\hline
Scenario & Excess at 99 & Excess at 198 & Deficit 100--149 & Hist.\ multi-filing & Hist.\ mean units \\
\hline
Unweighted & 11.4 & 2.9 & 7.9 & 6.7\% & 173 \\
Production & 11.4 & 2.9 & 7.5 & 5.3\% & 149 \\
Drop all flagged parents & 11.8 & 3.2 & 8.2 & 4.3\% & 146 \\
Earlier-parcel areas & 11.5 & 2.9 & 5.9 & 4.2\% & 124 \\
Flagged land halved & 11.4 & 2.9 & 7.5 & 5.7\% & 154 \\
Flagged land doubled & 11.4 & 2.9 & 7.5 & 5.0\% & 146 \\
\hline
\end{tabular}
\end{center}

Excesses and deficits are percentage points of the normalized parent-size
distribution; the post-period multi-filing share is 16.8 percent and
the post-period mean 137 units. The bootstrap percentile intervals
for the production benchmark are 8.0--14.9 points for the excess at
99 and 3.3--12.4 for the 100--149 deficit.

\paragraph{What it shows.} The excess at 99 and at 198 is the same under every
land definition, including no weighting at all; there are no 198-unit
historical parents in any scenario. Misstating every flagged parent's land by
a factor of two moves the excess at 99 by less than 0.02 points and
mean size by about five units, so the remaining per-parent queue
cannot change the headline.

The earlier-parcel scenario is different. It lowers the historical mean from
149 to 124 units, below the post-period mean, and the
100--149 deficit from 7.5 to 5.9 points. That rule enlarges
194 of 595 historical sites but only 33 of 310 post-period
sites. The reason is timing, not measurement quality: developers typically
file under the one tax lot that survives a later merger, and in 91
percent of the enlarged historical cases DOF recorded that merger after the
filing, typically about a year later. 589 Fulton Street is an example: filed in
August 2021 on lot 14, a 279-square-foot sliver in the January 2021 map, with
the nine-lot merger recorded in June 2022; the site is 27,542 square feet.
Post-period filings from 2025--2026 have mostly not yet had their mergers
recorded. Applying the rule therefore makes historical sites look larger than
comparable recent sites, and the weights move away from large historical
projects. The production measure, the filing lot in the pre-filing map, uses
the same method in both periods, and the smallest quarter of sites by land per
unit looks alike across periods (below 52 and 53 square feet
per unit); that does not establish that its errors balance.

\paragraph{Where this leaves the review.} The per-parent research batches were
not continued, and the September 22 cleanup retired their download rules and
memos (tag \texttt{pre-cleanup-2026-09-22} keeps them). Whether any case-level
review resumes is Jacob's decision. Adopted decisions remain in
\path{parent_opportunities_manual}; the automatic screen and candidate areas
remain in \path{audit_parent_site_boundaries}. Claude recommends treating the
open question as a rule rather than a list of cases: measure site land the
same way in both periods, for example by counting a DOF merger only if it was
recorded within a fixed window after filing, and add a check for implausible
land per unit. It matters for the size and housing-forgone questions, not for
the bunching.

\textit{Reproduction note.} Root \texttt{make land-sensitivity} runs
\path{tasks/audits/audit_land_measurement_sensitivity}, which reads the
canonical parent panel, the boundary screen and the saved benchmark weights,
and checks that its production scenario reproduces those weights. The merger
timing comes from \path{parent_site_scope.csv} and the DOF transaction dates in
\path{audit_parent_site_boundaries/output/dof_events.parquet}.
